\section{Residual-budget contracts}
\label{sec:contract}

The first half of the question asked which generator information explains the numerical gain,
and the answer was that block ownership does not. This half asks which generator information
can be turned into a guarantee. Ownership is one kind of generator information; a tolerance
stated in application units is another. A residual budget says how much violation of an original constraint is
acceptable. It must be supplied by the application or declared as a numerical test
parameter; it cannot be inferred from a block label or a zero right-hand side. Scaling rows
so that solver tolerances correspond to negligible application-unit violations is
established practice \cite{gurobiScaling,gurobiExternalScaling,idaesScaling}. The results
below make its combination with coefficient limits exact: they characterize every admissible
factor, certify when none exists, select the centered one, and identify what row scaling
cannot control. Individually the steps are elementary algebra; what the section contributes is
that the contract is complete, in the sense that a factor is admissible exactly when the
conditions hold, and that each consequence is stated in the units the application uses.

For a nonconstant row $a^Tx\le b$, define the original residual $v(x)=(a^Tx-b)_+$ and a
budget $\delta>0$. Let $\epsilon>0$ bound the residual of the \emph{exported} row. This is a
hypothesis about the returned point rather than a consequence of a solver parameter:
presolve, internal scaling and finite precision intervene between the two. For a positive
row factor $\alpha$,
\begin{equation}\label{eq:residual-transfer}
 (\alpha a^Tx-\alpha b)_+=\alpha v(x)
\end{equation}
in exact arithmetic, so an exported residual at most $\epsilon$ implies $v(x)\le\delta$
whenever $\alpha\ge\epsilon/\delta$. Absolute residuals give the same identity for
equalities. Write $m=\min_{a_j\ne0}|a_j|$ and $M=\max_j|a_j|$, and let $0<\ell\le U$ limit
the exported coefficient magnitudes. An additional RHS limit $U_b$ intersects the upper
endpoint below with $U_b/|b|$ for $b\ne0$.

\begin{proposition}[Admissible factors and constrained centering]
\label{prop:budget}
The factors that place every nonzero coefficient in $[\ell,U]$ and satisfy
$\epsilon/\alpha\le\delta$ form exactly the interval
\begin{equation}\label{eq:budget-interval}
 \mathcal I=[L,H],\qquad
 L=\max\{\epsilon/\delta,\ell/m\},\qquad H=U/M.
\end{equation}
The interval is nonempty if and only if
\begin{equation}\label{eq:budget-feasibility}
 M/m\le U/\ell\quad\text{and}\quad \delta\ge\epsilon M/U.
\end{equation}
When it is nonempty, the unique minimizer of
\begin{equation}\label{eq:budget-objective}
 \min_{\alpha\in\mathcal I}\;\max_{a_j\ne0}|\log(\alpha|a_j|)|
\end{equation}
is
\begin{equation}\label{eq:budget-optimum}
 \alpha^*=\min\{H,\max\{L,(mM)^{-1/2}\}\}.
\end{equation}
\end{proposition}

\begin{proof}
The coefficient bounds are equivalent to $\alpha m\ge\ell$ and $\alpha M\le U$.
Together with $\alpha\delta\ge\epsilon$ they give \eqref{eq:budget-interval}.
The condition $L\le H$ is equivalent to the two inequalities in
\eqref{eq:budget-feasibility}. To solve the objective, put $t=\log\alpha$,
$q=(\log m+\log M)/2$ and $h=(\log M-\log m)/2$. Since the extreme absolute
logarithms occur at an endpoint,
\[
 \max_j|t+\log|a_j||=\max\{|t+\log m|,|t+\log M|\}=h+|t+q|.
\]
The unique minimizer on $[\log L,\log H]$ is the projection of $-q$ onto that
interval. Exponentiating gives \eqref{eq:budget-optimum}.
\end{proof}

\paragraph{Incompatibility certificate.}
When $L>H$, \eqref{eq:budget-feasibility} names the violated requirement. A within-row
spread $M/m>U/\ell$ cannot be repaired by any row factor; a budget $\delta<\epsilon M/U$
conflicts with the coefficient ceiling even for a perfectly balanced row. In both cases the
certificate concerns the numerical specification, not the optimization model: exporting the
unscaled row does not meet the contract, and resolving the conflict requires revising units,
limits or the formulation. The factor $\alpha^*$ optimizes coefficient centering subject to
the contract; it is not a condition-number minimizer.

\begin{corollary}[Invariance under consistent re-emission]
\label{cor:budget-invariance}
Replace $(a,b,\delta)$ by $(qa,qb,q\delta)$ for any $q>0$, keeping
$\epsilon,\ell,U$ fixed. The interval becomes $\mathcal I/q$, abstention is
unchanged, and the selected factor becomes $\alpha^*/q$. The exported row and the
budget-normalized original residual $v(x)/\delta$ are therefore unchanged.
\end{corollary}
\begin{proof}
Both coefficient extremes and the original budget acquire the factor $q$.
Each term defining $L$, $H$ and $(mM)^{-1/2}$ is divided by $q$.
Equation~\eqref{eq:residual-transfer} gives the residual statement.
\end{proof}

A budget is therefore information that survives preprocessing: generators that emit the same
constraint in different units, each with the budget expressed in its own units, obtain the
same exported row and the same contract. The invariance is lost if a row is re-emitted
without transforming its budget, and a fixed cutoff on the factor itself does not have it.
Restricting factors to powers of two \cite{berthold2021scaling} retains it for power-of-two
re-emissions, subject to the representability conditions of \cref{prop:dyadic}.

\subsection{The integer-rounding term}
Let $J$ be the integer columns and let $\bar x$ be obtained by rounding only these
coordinates of $x$. Suppose $|\bar x_j-x_j|\le\eta_j$ for $j\in J$ and put
$E=\sum_{j\in J}|a_j|\eta_j$. The triangle inequality and \eqref{eq:residual-transfer} give
\begin{equation}\label{eq:rounding-budget}
 v(\bar x)\le \epsilon/\alpha+E,
\end{equation}
also for equalities with absolute residuals. Changing the row factor reduces the
exported-residual term but leaves the rounding term, which is measured in original units,
unchanged. Normalizing a big-$M$ row therefore cannot remove the effect of an integrality
tolerance on that row.

For a contract on the rounded assignment, \cref{prop:budget} applies with $\delta$ replaced by
$\delta-E$ when $\delta>E$; if $\delta\le E$, no positive factor certifies it. The bound is
a worst case: actual rounding errors may be smaller or cancel, and other rows can exclude
the worst case, so an empty interval shows that the stated information is insufficient
rather than that returned points will violate the budget. Under re-emission $E$ becomes $qE$,
so invariance is retained. A tighter integrality tolerance or a reformulation can reduce $E$;
row scaling cannot.

\subsection{Exact export on a dyadic grid}
A factor in the real admissible interval does not guarantee an exactly proportional exported
row: the products $\operatorname{fl}(\alpha a_j)$ and $\operatorname{fl}(\alpha b)$ are rounded
independently, and rounding a computed endpoint to nearest can leave the closed interval. Our
arbitrary-factor implementation forms the endpoints as exact rational numbers in the supplied
binary64 inputs, rounds them \emph{inward}, and abstains if no positive binary64 factor
remains. This certifies admissibility of the factor, not exactness of the products.

Powers of two avoid multiplication rounding. The following intersection also accounts for
subnormal numbers and for the RHS. Assume finite binary64 input data and IEEE arithmetic with
gradual underflow \cite{ieee2019float}. Write each nonzero datum $z\in\{a_1,\ldots,a_n,b\}$
uniquely as $z=\operatorname{sign}(z)N_z2^{p_z}$ with $N_z$ a positive odd integer, and put
$q_z=\lfloor\log_2|z|\rfloor$. Let $[L,H]$ be the desired real factor interval, including any
rounding allowance from \eqref{eq:rounding-budget}.

\begin{proposition}[Exact dyadic admissibility]
\label{prop:dyadic}
A factor $\alpha=2^k$ is positive and representable in binary64, belongs to
$[L,H]$, and multiplies every nonzero coefficient and the RHS exactly into a
finite binary64 value if and only if the integer $k$ satisfies
\begin{align*}
 K_-&=\max\left\{-1074,\lceil\log_2L\rceil,
                    \max_z(-1074-p_z)\right\}\le k,\\
 k&\le K_+=\min\left\{1023,\lfloor\log_2H\rfloor,
                    \min_z(1023-q_z)\right\}.
\end{align*}
If $K_->K_+$ there is no such factor. Otherwise an exponent nearest to
$-\tfrac12\log_2(mM)$ among the integers in $[K_-,K_+]$ minimizes the
logarithmic coefficient objective \eqref{eq:budget-objective} on this exact-export
set; a tie may be resolved by choosing the larger factor.
\end{proposition}
\begin{proof}
The factor itself requires $-1074\le k\le1023$. Since the datum's odd
significand is unchanged, the product $N_z2^{p_z+k}$ lies on the subnormal grid
exactly when $p_z+k\ge-1074$; its highest bit stays below overflow exactly when
$q_z+k\le1023$. The original significand has at most 53 bits, so these two
conditions are also sufficient. Intersecting them over the row data and with
$L\le2^k\le H$ gives the interval of integer exponents. The objective is
$\tfrac12\log(M/m)+|k\log2+\tfrac12\log(mM)|$, proving the nearest-exponent
claim.
\end{proof}

The implementation obtains the endpoint exponents from integer ratios and bit lengths,
compares candidate objectives exactly as rational numbers, and checks the rounded products
against exact rational products. Writing the row to a text file additionally requires a
round-trip decimal serializer, since fixed decimal places can discard small values even when
the multiplication is exact.

\paragraph{Verification after solving.}
The contract is conditional on the exported residual, so returned points are checked against
the original rows, bounds and integrality restrictions. On the stress instances of
\cref{sec:experiments} the check is exact: row residuals are rational expressions in the stored
binary64 data, and optima are obtained by enumeration. On the dispatch models the archived
checker uses floating-point residuals, which are diagnostics; certifying branch-and-bound
results is a stronger task \cite{cheung2017certificates}.
